\documentclass[a4paper,11pt]{article}

% Encoding and language
\usepackage[utf8]{inputenc}
\usepackage[T1]{fontenc}
\usepackage[english]{babel} % change to dutch if needed

% Page layout
\usepackage[a4paper,margin=2.5cm]{geometry}
\usepackage{setspace}
\usepackage{parskip}

% Math and symbols
\usepackage{amsmath,amssymb}

% Figures and tables
\usepackage{graphicx}
\usepackage{float}
\usepackage{subcaption}
\usepackage{booktabs}
\usepackage{array}

% Links
\usepackage[hidelinks]{hyperref}

% Better lists
\usepackage{enumitem}

% Optional: nicer captions
\usepackage[font=small,labelfont=bf]{caption}

\title{\textbf{Simulation of Process Scheduling Algorithms}\\[0.5em]
Operating Systems 1 -- Subproject 1}
\author{
    Student 1 Name \and
    Student 2 Name
}
\date{\today}

\begin{document}

\maketitle

\vspace{-1em}
\begin{abstract}
This report presents a simulation-based evaluation of several CPU scheduling
algorithms: FCFS, RR, SPN, SRT, HRRN, and MLFB. We compare their performance
on the provided datasets using average turnaround time (TAT), normalized turnaround
time (nTAT), waiting time ($T_w$), and percentile-based graphs. We also motivate
our chosen MLFB parameters and compare the observed behaviour with the theoretical
results from the lecture material.
\end{abstract}

\section{Experimental Setup}

This section briefly describes the simulator and the experiments.

\subsection{Implemented Scheduling Algorithms}
We implemented the following scheduling algorithms:
\begin{itemize}[leftmargin=1.5em]
    \item First Come, First Serve (FCFS)
    \item Round Robin (RR) with time slices 2, 4, and 8
    \item Shortest Process Next (SPN)
    \item Shortest Remaining Time (SRT)
    \item Highest Response Ratio Next (HRRN)
    \item Multi-Level Feedback (MLFB) with 5 levels
\end{itemize}

\subsection{Datasets}
We used the provided XML datasets containing processes with:
\begin{itemize}[leftmargin=1.5em]
    \item process ID,
    \item arrival time,
    \item service time.
\end{itemize}

The experiments were performed on the following datasets:
\begin{itemize}[leftmargin=1.5em]
    \item Dataset A: 10000 processes
    \item Dataset B: 20000 processes
    \item Dataset C: 50000 processes
\end{itemize}

\subsection{Measured Metrics}
For each scheduling algorithm and each dataset, we computed:
\begin{itemize}[leftmargin=1.5em]
    \item average turnaround time (TAT),
    \item average normalized turnaround time (nTAT),
    \item average waiting time ($T_w$).
\end{itemize}

In addition, for two selected datasets, we generated:
\begin{itemize}[leftmargin=1.5em]
    \item one $nTAT(T_s)$ graph,
    \item one $T_w(T_s)$ graph,
\end{itemize}
where each point represents the arithmetic mean over one percentile of service times.

\section{Design and Implementation}

\subsection{Simulator Structure}
Briefly describe the structure of your simulator:
\begin{itemize}[leftmargin=1.5em]
    \item how XML input is parsed,
    \item how processes are represented,
    \item how time progression is simulated,
    \item how scheduling decisions are made,
    \item how statistics are collected.
\end{itemize}

\subsection{Assumptions and Design Choices}
Mention any relevant implementation decisions, for example:
\begin{itemize}[leftmargin=1.5em]
    \item tie-breaking rules,
    \item ready queue data structures,
    \item treatment of arrivals at the same time,
    \item exact interpretation of one jiffy and rescheduling moments.
\end{itemize}

\section{MLFB Configuration}

We used an MLFB scheduler with 5 levels. The chosen parameters are:

\begin{table}[H]
    \centering
    \begin{tabular}{@{}ll@{}}
        \toprule
        Parameter & Value / Description \\
        \midrule
        Number of levels & 5 \\
        Time quantum per level & e.g.\ 1, 2, 4, 8, 16 \\
        New job insertion level & e.g.\ highest priority queue \\
        Demotion rule & e.g.\ after full quantum usage \\
        Promotion / aging rule & e.g.\ none or periodic boost \\
        Priority boost interval & e.g.\ every 50 jiffies / not used \\
        \bottomrule
    \end{tabular}
    \caption{Chosen MLFB parameters.}
    \label{tab:mlfb}
\end{table}

\subsection{Motivation}
Explain why these parameters were chosen. For example, discuss the trade-off between:
\begin{itemize}[leftmargin=1.5em]
    \item responsiveness for short jobs,
    \item fairness,
    \item starvation avoidance,
    \item approximation of SRT-like behaviour.
\end{itemize}

\section{Results}

\subsection{Global Measurements}
Insert the global averages for each algorithm and dataset.

\begin{table}[H]
    \centering
    \begin{tabular}{@{}lllll@{}}
        \toprule
        Dataset & Algorithm & Avg.\ TAT & Avg.\ nTAT & Avg.\ $T_w$ \\
        \midrule
        10000 & FCFS   &  &  &  \\
        10000 & RR-2   &  &  &  \\
        10000 & RR-4   &  &  &  \\
        10000 & RR-8   &  &  &  \\
        10000 & SPN    &  &  &  \\
        10000 & SRT    &  &  &  \\
        10000 & HRRN   &  &  &  \\
        10000 & MLFB   &  &  &  \\
        \midrule
        20000 & FCFS   &  &  &  \\
        20000 & RR-2   &  &  &  \\
        20000 & RR-4   &  &  &  \\
        20000 & RR-8   &  &  &  \\
        20000 & SPN    &  &  &  \\
        20000 & SRT    &  &  &  \\
        20000 & HRRN   &  &  &  \\
        20000 & MLFB   &  &  &  \\
        \midrule
        50000 & FCFS   &  &  &  \\
        50000 & RR-2   &  &  &  \\
        50000 & RR-4   &  &  &  \\
        50000 & RR-8   &  &  &  \\
        50000 & SPN    &  &  &  \\
        50000 & SRT    &  &  &  \\
        50000 & HRRN   &  &  &  \\
        50000 & MLFB   &  &  &  \\
        \bottomrule
    \end{tabular}
    \caption{Global measurements for all algorithms and datasets.}
    \label{tab:global-results}
\end{table}

\subsection{Graphs for Selected Datasets}

\subsubsection{Dataset X}
\begin{figure}[H]
    \centering
    \includegraphics[width=0.8\textwidth]{figures/datasetX_ntat.png}
    \caption{$nTAT(T_s)$ for Dataset X.}
    \label{fig:datasetX-ntat}
\end{figure}

\begin{figure}[H]
    \centering
    \includegraphics[width=0.8\textwidth]{figures/datasetX_tw.png}
    \caption{$T_w(T_s)$ for Dataset X.}
    \label{fig:datasetX-tw}
\end{figure}

\subsubsection{Dataset Y}
\begin{figure}[H]
    \centering
    \includegraphics[width=0.8\textwidth]{figures/datasetY_ntat.png}
    \caption{$nTAT(T_s)$ for Dataset Y.}
    \label{fig:datasetY-ntat}
\end{figure}

\begin{figure}[H]
    \centering
    \includegraphics[width=0.8\textwidth]{figures/datasetY_tw.png}
    \caption{$T_w(T_s)$ for Dataset Y.}
    \label{fig:datasetY-tw}
\end{figure}

\section{Evaluation and Comparison}

\subsection{Overall Comparison}
Discuss the main trends in the global measurements. For example:
\begin{itemize}[leftmargin=1.5em]
    \item Which algorithms minimize average TAT and nTAT?
    \item Which algorithms favour short jobs most strongly?
    \item How does RR change as the quantum increases?
    \item How competitive is MLFB compared to SPN/SRT?
\end{itemize}

\subsection{Behaviour Across Service-Time Percentiles}
Discuss the four graphs:
\begin{itemize}[leftmargin=1.5em]
    \item performance on short processes,
    \item performance on long processes,
    \item fairness,
    \item whether some algorithms disproportionately penalize certain job sizes.
\end{itemize}

\subsection{Round Robin Analysis}
Comment on the effect of time slice choice:
\begin{itemize}[leftmargin=1.5em]
    \item small quantum $\rightarrow$ more responsiveness,
    \item large quantum $\rightarrow$ behaviour closer to FCFS,
    \item trade-off between short-job responsiveness and waiting time for long jobs.
\end{itemize}

\subsection{MLFB Analysis}
Evaluate whether your MLFB parameterization achieves its intended goal. Explain:
\begin{itemize}[leftmargin=1.5em]
    \item which scheduling behaviour it resembles most,
    \item where it performs well,
    \item where it performs poorly,
    \item whether priority boosting or demotion policy had a visible effect.
\end{itemize}

\section{Comparison with Theory}

Compare your results with the theoretical behaviour presented in the lectures.

Points to address:
\begin{itemize}[leftmargin=1.5em]
    \item whether SPN and SRT outperform FCFS and RR for short jobs,
    \item whether HRRN balances efficiency and fairness as expected,
    \item whether MLFB approximates short-job preference,
    \item whether the observed trends match the lecture material qualitatively.
\end{itemize}

If some results differ from theory, acknowledge this clearly and discuss possible causes, such as:
\begin{itemize}[leftmargin=1.5em]
    \item implementation errors,
    \item tie-breaking effects,
    \item different MLFB parameter settings,
    \item plotting or averaging choices,
    \item simulation assumptions.
\end{itemize}

\section{Conclusion}

Summarize the most important findings in a few paragraphs. Keep this section focused:
\begin{itemize}[leftmargin=1.5em]
    \item best-performing algorithms overall,
    \item trade-offs between fairness and efficiency,
    \item impact of RR quantum,
    \item lessons learned from MLFB.
\end{itemize}

% Optional
\section*{Appendix (Optional)}
Use this only if really necessary, for example for a small extra table or a short note
on implementation details. Do not include source code.

\end{document}